% GENERATED FILE. Edit canonical lessons, then run npm run generate:books.
% canonical-source-hash: fnv1a64:a1cef325c44c0980
% canonical-lessons: TE-C261-sahajamaina, TE-C261-cikkani, TE-C261-jiddu, TE-C261-jigata, TE-C261-garuku

\chapter{Natural, Thick, Greasy, Sticky, Rough}
\label{ch:te-natural-thick-greasy-sticky-rough}
\hlchaptermodality{261}

\begin{chapteropening}
\textbf{By the end of this chapter:} \emph{I can say natural, thick, greasy, sticky and rough.}

Complete the last lesson of chapter 261: I can say natural, thick, greasy, sticky and rough.
\end{chapteropening}

\section[\texorpdfstring{\te{సహజమైన} (sahajamaina)}{సహజమైన (sahajamaina)}]{\te{సహజమైన} (sahajamaina) — natural}

\label{lesson:TE-C261-sahajamaina}



\begin{quote}
Before the new one: say the Telugu for strange, then the Telugu for kind.
\end{quote}

\subsection*{You'll want to know: \te{సహజమైన}}
\textbf{\te{సహజమైన}} — \emph{sahajamaina} — \textquotedblleft{}natural\textquotedblright{}.

\begin{grammarlens}[title={What you've built}]
One more word for this chapter.
\end{grammarlens}

\subsection*{Your turn}
\begin{itemize}
\raggedright
  \item \emph{Say it:} \emph{sahajamaina}
  \item \emph{Say it:} \emph{sahajamaina}, once more
  \item \emph{Say it:} say \emph{sahajamaina}
  \item \emph{Recall:} say the Telugu for strange, then the Telugu for kind, then say \emph{sahajamaina} again
\end{itemize}

\begin{tcolorbox}[breakable,colback=teal!4,colframe=teal!35!black,title={Before you move on}]
What does \te{సహజమైన} mean? (\textquotedblleft{}Natural\textquotedblright{}.) Say it once more.
\end{tcolorbox}
\section[\texorpdfstring{\te{చిక్కని} (cikkani)}{చిక్కని (cikkani)}]{\te{చిక్కని} (cikkani) — thick (of a liquid)}

\label{lesson:TE-C261-cikkani}



\begin{quote}
Before the new one: say the Telugu for kind, then the Telugu for natural.
\end{quote}

\subsection*{You'll want to know: \te{చిక్కని}}
\textbf{\te{చిక్కని}} — \emph{cikkani} — \textquotedblleft{}thick (of a liquid)\textquotedblright{}.

\begin{grammarlens}[title={What you've built}]
One more word for this chapter.
\end{grammarlens}

\subsection*{Your turn}
\begin{itemize}
\raggedright
  \item \emph{Say it:} \emph{cikkani}
  \item \emph{Say it:} \emph{cikkani}, once more
  \item \emph{Say it:} say \emph{cikkani}
  \item \emph{Recall:} say the Telugu for kind, then the Telugu for natural, then say \emph{cikkani} again
\end{itemize}

\begin{tcolorbox}[breakable,colback=teal!4,colframe=teal!35!black,title={Before you move on}]
What does \te{చిక్కని} mean? (\textquotedblleft{}Thick (of a liquid)\textquotedblright{}.) Say it once more.
\end{tcolorbox}
\section[\texorpdfstring{\te{జిడ్డు} (jiḍḍu)}{జిడ్డు (jiḍḍu)}]{\te{జిడ్డు} (jiḍḍu) — greasy, oily}

\label{lesson:TE-C261-jiddu}



\begin{quote}
Before the new one: say the Telugu for natural, then the Telugu for thick.
\end{quote}

\subsection*{You'll want to know: \te{జిడ్డు}}
\textbf{\te{జిడ్డు}} — \emph{jiḍḍu} — \textquotedblleft{}greasy, oily\textquotedblright{}.

\begin{grammarlens}[title={What you've built}]
One more word for this chapter.
\end{grammarlens}

\subsection*{Your turn}
\begin{itemize}
\raggedright
  \item \emph{Say it:} \emph{jiḍḍu}
  \item \emph{Say it:} \emph{jiḍḍu}, once more
  \item \emph{Say it:} say \emph{jiḍḍu}
  \item \emph{Recall:} say the Telugu for natural, then the Telugu for thick, then say \emph{jiḍḍu} again
\end{itemize}

\begin{tcolorbox}[breakable,colback=teal!4,colframe=teal!35!black,title={Before you move on}]
What does \te{జిడ్డు} mean? (\textquotedblleft{}Greasy, oily\textquotedblright{}.) Say it once more.
\end{tcolorbox}
\section[\texorpdfstring{\te{జిగట} (jigaṭa)}{జిగట (jigaṭa)}]{\te{జిగట} (jigaṭa) — sticky}

\label{lesson:TE-C261-jigata}



\begin{quote}
Before the new one: say the Telugu for thick, then the Telugu for greasy.
\end{quote}

\subsection*{You'll want to know: \te{జిగట}}
\textbf{\te{జిగట}} — \emph{jigaṭa} — \textquotedblleft{}sticky\textquotedblright{}.

\begin{grammarlens}[title={What you've built}]
One more word for this chapter.
\end{grammarlens}

\subsection*{Your turn}
\begin{itemize}
\raggedright
  \item \emph{Say it:} \emph{jigaṭa}
  \item \emph{Say it:} \emph{jigaṭa}, once more
  \item \emph{Say it:} say \emph{jigaṭa}
  \item \emph{Recall:} say the Telugu for thick, then the Telugu for greasy, then say \emph{jigaṭa} again
\end{itemize}

\begin{tcolorbox}[breakable,colback=teal!4,colframe=teal!35!black,title={Before you move on}]
What does \te{జిగట} mean? (\textquotedblleft{}Sticky\textquotedblright{}.) Say it once more.
\end{tcolorbox}
\section[\texorpdfstring{\te{గరుకు} (garuku)}{గరుకు (garuku)}]{\te{గరుకు} (garuku) — rough (to the touch)}

\label{lesson:TE-C261-garuku}



\begin{quote}
Before the new one: say the Telugu for greasy, then the Telugu for sticky.
\end{quote}

\subsection*{You'll want to know: \te{గరుకు}}
\textbf{\te{గరుకు}} — \emph{garuku} — \textquotedblleft{}rough (to the touch)\textquotedblright{}.

\begin{grammarlens}[title={What you've built}]
One more word for this chapter.
\end{grammarlens}

\subsection*{Your turn}
\begin{itemize}
\raggedright
  \item \emph{Say it:} \emph{garuku}
  \item \emph{Say it:} \emph{garuku}, once more
  \item \emph{Say it:} say \emph{garuku}
  \item \emph{Recall:} say the Telugu for greasy, then the Telugu for sticky, then say \emph{garuku} again
\end{itemize}

\begin{tcolorbox}[breakable,colback=teal!4,colframe=teal!35!black,title={Before you move on}]
What does \te{గరుకు} mean? (\textquotedblleft{}Rough (to the touch)\textquotedblright{}.) Say it once more.
\end{tcolorbox}
